\begin{titlepage}
  \centering
  \vspace*{30mm}
  {\Huge\bfseries Programming the machine\par}
  \vspace{8mm}
  {\Large Assembly programming with MARS\par}
  \vspace{18mm}
  {\Large\MARSAuthor\par}
  \vfill
  {\MARSEdition\par}
\end{titlepage}

\chapter*{Using this book}
\markboth{Using this book}{Using this book}
\addcontentsline{toc}{chapter}{Using this book}
\label{revision:using-learning-book}

This book assumes that you have written programs with variables, expressions,
conditionals, loops and functions. It assumes no assembly experience. The opening
chapters build a working model through small complete programs. Later chapters
use that model to reason about arrays, interfaces, vector operations and files.
You do not need to memorize an instruction catalog before beginning.

Read the explanation and trace a concrete input before running the source. At a
prediction, write down the value or state you expect and why. After execution,
compare that reasoning with the observed state. When they disagree, find the
first instruction or assumption that accounts for the difference. A correct final
output is useful evidence, but an explanation of its intermediate state is what
lets you change the program independently.

The main book completes smaller programs within its narrative. Complete public
sources for longer examples are in \path{revision/labs/files}. The accompanying
\emph{MARS laboratory companion} has the same numbered chapter sequence. It
provides the detailed contracts, observations, diagnosis exercises and independent
tasks associated with those chapters. A partial listing in this book has a
stated reading purpose; use its named complete source when executing it.

Follow the foundational sequence before the vector and file chapters. Elective
companion assignments offer additional practice with devices, data structures,
text and arithmetic. Completing the cryptography or interpreter assignments is
not a prerequisite for the audio, image or video sequence. Each chapter and lab
states the concepts it expects you to know.

The examples use the MARS MIPS64 Release 6 environment and its educational syscall
interface. Use the matching course plugin with IntelliJ IDEA 2026.3 and JDK 25.
Instruction behavior, data layout and interface details are explained where the
programs need them. Keep the source, assignment and checks from the same course
edition when comparing results.

The contents lead to chapters. The glossary collects terminology for reference;
it does not replace the first explanations. The subject index points to the
chapters that develop each topic.
